
\chapter{Implementation Details}
\label{app:implementation}

This appendix describes the per-step computation graph of the encoders we
tested, as well as the general CMVF training loop. All four encoder types -- GRU, LSTM, Transformer, and Mamba -- share the same input pipeline, the same bounded readout, and the
same mechanistic step; only the sequence model in the middle changes. We
therefore present the pipeline once and then list the per-encoder differences.

% ---------------------------------------------------------------------
\section{Per-step pipeline}
\label{app:pipeline}

At observation step $k$ each encoder consumes the current control vector
$u_k \in \mathbb{R}^{U}$ and the previous (predicted or teacher-forced) state
$\hat{y}_{k-1} \in \mathbb{R}^{P}$, and emits a positive parameter vector
$\theta_k \in \mathbb{R}^{m}_{>0}$ that is then fed to the mechanistic step
(Sec.~\ref{app:step}). The pipeline is
\begin{align}
\tilde{u}_k &= \phi_u(u_k), \qquad \tilde{y}_{k-1} = \phi_y(\hat{y}_{k-1}), \\
\mathrm{feat}_k &= \big[\,\tilde{u}_k\,;\,\tilde{y}_{k-1}\,\big], \\
x_k &= \mathrm{Lift}(\mathrm{feat}_k), \\
z_k &= \mathrm{Encoder}\big(x_{1:k}\big), \\
r_k &= \mathrm{Head}(z_k), \qquad \theta_k = \gamma(r_k).
\end{align}

where $\phi_u, \phi_y$ are scalar feature transforms; all synthetic data scenarios (glycolysis, MOF, and enzyme) have no feature transformations ($\phi_u(u_k),\phi_u(\hat{y}_{k-1} = \mathbb{I}$). For the IVTT runs we use
$\phi_u(u) = \sqrt{u}$ and
$\phi_y(y) = \max(\sqrt{y}, 1)$.

Lift and encoder dimensions, along with the other key training hyperparameters for the synthetic scenarios are summarized in Table \ref{tab:synthetic_hparams}. 

 DNA is collapsed to a single input at the first time step for all trajectories, and excluded from the encoder input. Furthermore, in this scenario, only the control amounts $u_k$ are known for the inputs in the reaction. The only $y_k$ trajectories observed are the two final  species (mature mRNA
\texttt{mm} and mature protein \texttt{pm}) feed back into the encoder. Similarly, for the IVTT runs, training hyperparameters are found can be found in Table \ref{tab:ivtt-hparams}. 

\paragraph{Gamma bounds on learned parameters.} In the pipeline, a single linear head produces a
logit $r_k \in \mathbb{R}^{m}$. We map this into a per-parameter interval $[\theta^{\mathrm{lo}}, \theta^{\mathrm{hi}}]$ to incorporate an additional prior regarding the magnitude of the learned $\theta$ that also serves to stabilize training. This is done with a log-sigmoid bound:
\begin{equation}
\gamma(r) \;=\; \theta^{\mathrm{lo}} \,\exp\!\Big(\log\!\big(\theta^{\mathrm{hi}}/\theta^{\mathrm{lo}}\big)\,\sigma(r/\tau)\Big),
\label{eq:loggamma}
\end{equation}
where $\sigma$ is the elementwise sigmoid and $\tau\!\geq\!1$ a temperature.
At $r\!=\!0$ this maps to $\sqrt{\theta^{\mathrm{lo}}\theta^{\mathrm{hi}}}$,
which is essential because per-parameter bounds span different orders of
magnitude (e.g.\ $k_{\mathrm{nuc},A}\in[10^{-6},10^{-1}]$). A linear sigmoid would initialize at the arithmetic mean and blows up the ODE on the first
forward pass. We use $\tau\!=\!2.3$ for all encoders.

\paragraph{Teacher forcing.} During training, every \texttt{tf\_every}-th step
the previous-state input $\hat{y}_{k-1}$ is replaced by the ground-truth
$y_{k-1}$ on the observed channels only. Teacher forcing is disabled after
epoch \texttt{tf\_drop\_epoch} so the final training phase is fully
free-running.

% ---------------------------------------------------------------------
\section{Encoder variants}
\label{app:encoders}

Each encoder takes $\mathrm{feat}_k$ and returns a $d_{\mathrm{hidden}}$-d
context $z_k$. The same four variants are used across all scenarios:
on the synthetic benchmarks (MOF, Glycolysis, Enzyme) we sweep the four
encoders against a clean simulation oracle, and on the real IVTT data we
re-use the identical encoder code with only width/depth and the input
transforms changed (cf.\ Tables~\ref{tab:synthetic_hparams}
and~\ref{tab:ivtt-hparams}). Only the parts that differ from
Sec.~\ref{app:pipeline} are shown below.

\paragraph{GRU.} $L$ GRU layers with dropout applied to every layer's
output (including the last) during training, and orthogonal $W_{hh}$ init, Xavier $W_{ih}$ init (zero biases). 
\begin{equation}
h^{(\ell)}_k = \mathrm{GRUCell}^{(\ell)}\!\big(h^{(\ell-1)}_k,\, h^{(\ell)}_{k-1}\big), \quad \ell=1,\dots,L, \qquad z_k = h^{(L)}_k.
\end{equation}

\paragraph{LSTM.}  Standard PyTorch
\texttt{nn.LSTM} with $L$ stacked LSTM layers of width
$d_{\mathrm{hidden}}$. We use PyTorch's default initialisation.

\paragraph{Transformer.} Standard PyTorch \texttt{nn.TransformerEncoder} implementation, with a sliding window, where the transformer only attends to the last $W$ features. To account for this, we use a window-relative learned positional embedding:
\begin{align}
x_k &= \mathrm{Lift}(\mathrm{feat}_k) + \mathrm{PosEmbed}(k\!\!\mod\!\!W), \\
z_k &= \mathrm{TransformerEnc}_{\mathrm{causal}}\!\big(x_{\max(1,k-W+1):k}\big)_{[-1]}.
\end{align}
The number of heads is $d_{\mathrm{hidden}}/32$ and the feed-forward width is
\texttt{ff\_mult}$\,\cdot d_{\mathrm{hidden}}$. We detach features that scroll
out of the context window so memory is bounded at
$\mathcal{O}(K\,W)$ activations.

\paragraph{Mamba.} $L$ pre-norm Mamba blocks with residual
connections,
\begin{equation}
x^{(\ell)}_k = x^{(\ell-1)}_k + \mathrm{Mamba}^{(\ell)}\!\big(\mathrm{LayerNorm}(x^{(\ell-1)}_k)\big),
\end{equation}
and read out $z_k = x^{(L)}_k$. Because the closed-loop feedback
$\hat{y}_{k-1}\!\to\!x_k$ rules out a parallel scan, we run Mamba in
recurrent mode via \texttt{InferenceParams}.

% ---------------------------------------------------------------------
\subsection{Mechanistic step}
\label{app:step}

Given $\theta_k$ and the previous state $\hat{y}_{k-1}$, we advance the
state to $\hat{y}_k$ over the irregular timestep $\Delta t_k$ using one of
two paths, selected by the scaffold.

\paragraph{Numerical integration (synthetic scenarios + TXTL maturation).}
This is the default path used for the MOF, Glycolysis and Enzyme synthetic
benchmarks, and re-used for the
\texttt{txtl\_resource\_and\_maturation\_dna} scaffold. Control inputs enter
the state at the start of each interval as an instantaneous bolus, after
which the scaffold's right-hand side $f$ is integrated with one RK4 step:
\begin{align}
\hat{y}_k^{-} &= \hat{y}_{k-1} + J\,u_k, \\
\hat{y}_k     &= \mathrm{RK4}\!\big(\hat{y}_k^{-},\; f(\,\cdot\,;\theta_k),\; \Delta t_k\big).
\end{align}
Here $J\!\in\!\mathbb{R}^{P\times U}$ is the fixed control-to-state jump
matrix. State outputs are clamped at zero after each step for mass
conservation. We never differentiate through the integrator's substep loop
beyond the current observation interval, since teacher forcing /
\texttt{detach\_y\_prev} resets the gradient chain at every $k$.

\paragraph{Closed-form step (IVTT only).}  The IVTT scaffold
(\texttt{IvttAnalyticScaffold}) provides a closed-form solution of the
mature-mRNA / mature-protein ODE between observations, so we bypass the RK4
integrator entirely. A per-batch context (cumulative DNA) is precomputed
once and routed directly into the closed-form expression, and at every
step the scaffold's \texttt{analytic\_step} returns $\hat{y}_k$ directly as
a function of $(\hat{y}_{k-1}, \theta_k, \Delta t_k)$.

% ---------------------------------------------------------------------
\subsection{Loss}
\label{app:loss}

For all synthetic scenarios (MOF, Glycolysis, Enzyme) and for the TXTL
maturation runs we use a plain mean-squared error in $\log(1+\cdot)$ space,
averaged over the observed species $\mathcal{O}$ and over valid timesteps:
\begin{equation}
\mathcal{L}
\;=\;
\frac{1}{B\,K\,|\mathcal{O}|}
\sum_{b,k,j\in\mathcal{O}}
\Big(\log(1+\hat{y}^{(b)}_{k,j}) - \log(1+y^{(b)}_{k,j})\Big)^2 .
\label{eq:loss-base}
\end{equation}
The $\log(1+\cdot)$ wrap is essential because state magnitudes span several
orders of magnitude across species; a raw-scale MSE would let the largest
species dominate the gradient.

For the IVTT sweep we keep the same $\log(1+\cdot)$ space but replace the
uniform weighting in Eq.~\ref{eq:loss-base} with a three-term weighted MSE
that targets the two species we actually observe in the wet-lab data:
\begin{enumerate}
  \item a mature-mRNA term, weighted on every valid step;
  \item a mature-protein term, with second-half steps up-weighted by a
        factor of $N=3$;
  \item an endpoint term on the final mature-protein value.
\end{enumerate}
The remaining channels (resource, immature mRNA/protein, and proxies on
the emitted parameters) carry zero weight in the final aggregation --
they are computed for diagnostics only.

We optimise with AdamW, gradient clipping at $\|\nabla\|_2\!\leq\!1$, and
optionally a linear-warmup schedule (Transformer / Mamba). All models are
trained for 250 epochs on a fixed 0.125 / 0.125 test / validation split
hard-coded in \texttt{bob\_model/3 seed splits.py} so the same held-out
samples are scored across encoders.

% ---------------------------------------------------------------------
\subsection{Hyperparameters}
\label{app:hparams}

Tables~\ref{tab:synthetic_hparams} and~\ref{tab:ivtt-hparams} list the
hyperparameters used for the synthetic-scenarios sweep and for the IVTT
sweep, respectively. In both tables the settings that are identical across
all four encoders are folded into the caption, and the body shows only
per-encoder differences.

\begin{table}[h]
\centering
\caption{IVTT encoder ablation hyperparameters. The GRU and LSTM are
trained on \texttt{real\_ivtt\_exact\_full\_drop5000.npz} (a filtered
variant), while the Transformer and Mamba are trained on
\texttt{real\_ivtt\_exact\_full.npz}; both share the same fixed
$0.125\,/\,0.125$ test\,/\,validation split. All runs use 250 epochs,
optimiser AdamW with weight decay $1.8\!\times\!10^{-3}$, gradient clipping
at $1.0$, the per-parameter $\theta$ bounds defined by
\texttt{IvttAnalyticScaffold}, and the bounded readout in
Eq.~\ref{eq:loggamma}. Encoder feature mask: drop DNA from $u$, keep only
\texttt{mm} and \texttt{pm} from $y$. Loss: \texttt{loss\_fn\_ivtt\_mse}.}
\label{tab:ivtt-hparams}
\small
\begin{tabular}{lcccc}
\hline
                        & GRU                  & LSTM                 & Transformer                                 & Mamba \\
\hline
$d_{\mathrm{hidden}}$   & 400                  & 400                  & 128                                         & 400 \\
$L$ (layers)            & 2                    & 2                    & 2                                           & 2 \\
dropout                 & 0.2                  & 0.2                  & 0.1                                         & 0.2 \\
lift                    & none                 & $\to\!32$            & $\to\!32\!\to\!128$                         & $\to\!32\!\to\!400$ \\
$\phi_u$                & $\sqrt{\cdot}$       & $\sqrt{\cdot}$       & $\mathrm{cumsum}\!\circ\!\sqrt{\cdot}$      & $\sqrt{\cdot}$ \\
$\phi_y$                & $\max(\sqrt{\cdot},1)$ & $\max(\sqrt{\cdot},1)$ & $\max(\sqrt{\cdot},1)$                  & $\max(\sqrt{\cdot},1)$ \\
$\tau$ (Eq.~\ref{eq:loggamma}) & 2.3           & 1.0                  & 1.0                                         & 1.0 \\
batch size              & 100                  & 100                  & 32                                          & 100 \\
learning rate           & $2\!\times\!10^{-3}$ & $2\!\times\!10^{-3}$ & $2\!\times\!10^{-3}$                        & $2\!\times\!10^{-3}$ \\
warmup epochs           & 0                    & 0                    & 0                                           & 10 \\
TF interval / drop      & 200 / 100            & 200 / 100            & 200 / 100                                   & 200 / 100 \\
encoder-specific        & ---                  & ---                  & $W\!=\!180$, $\mathrm{ff}\!=\!2d$           & $d_{\mathrm{state}}\!=\!16$, exp.\ 2, $d_{\mathrm{conv}}\!=\!4$ \\
\hline
\end{tabular}
\end{table}

\begin{table}[h]
\centering
\caption{Synthetic-scenarios encoder ablation. The same four-encoder sweep
is run on each of the Enzyme, Glycolysis, MOF and Methane datasets. All
runs use AdamW with weight decay $8.5\!\times\!10^{-3}$, gradient clipping
at $1.0$, scalar $\theta$ bounds $[10^{-3}, 2]$, the bounded readout in
Eq.~\ref{eq:loggamma} with $\tau\!=\!1$, no input transforms
($\phi_u\!=\!\phi_y\!=\!\mathrm{id}$), 100 epochs of training with a
fixed $800/100/100$ train / val / test split, batch size $32$,
learning rate $1\!\times\!10^{-3}$, and $n_{\mathrm{substeps}}\!=\!1$ in the
RK4 integrator. Per-encoder rows below list only what differs.}
\label{tab:synthetic_hparams}
\small
\begin{tabular}{lcccc}
\hline
                     & GRU                     & LSTM                & Transformer                                & Mamba \\
\hline
$d_{\mathrm{hidden}}$ & 128                    & 128                 & 128                                        & 128 \\
$L$ (layers)         & 2                       & 2                   & 2                                          & 2 \\
dropout              & 0.1                     & 0.1                 & 0.0                                        & 0.1 \\
lift                 & $\to\!32$               & $\to\!32$           & $\to\!32\!\to\!128$                        & $\to\!32\!\to\!128$ \\
warmup epochs        & 0                       & 0                   & 10                                         & 0 \\
TF interval / drop   & 50 / 500                & 50 / 500            & 50 / 500                                   & 50 / 500 \\
encoder-specific     & ---                     & ---                 & $W\!=\!128$, $\mathrm{ff}\!=\!4d$          & $d_{\mathrm{state}}\!=\!16$, exp.\ 2, $d_{\mathrm{conv}}\!=\!4$ \\
\hline
\end{tabular}
\end{table}

The corresponding YAML configs are
\texttt{configs/ivtt\_gru.yaml}, \texttt{configs/ivtt\_lstm.yaml},
\texttt{configs/ivtt\_transformer.yaml}, and
\texttt{configs/ivtt\_mamba\_ssm.yaml}; running
\texttt{python last-layer-ode/train.py --config <yaml>} reproduces the
corresponding row of Table~\ref{tab:ivtt-hparams}.
